\section{Methodology}
\label{sec:method}

This section describes how the original repository implements SGNS and what we changed. Figure~\ref{fig:pipeline} gives the overall pipeline; Appendix~\ref{app:code} maps each concept to its location in the code.

\begin{figure}[t]
    \centering
    \resizebox{\linewidth}{!}{% Training and evaluation pipeline with the code location of each stage (shared by report and slides).
% Requires \usetikzlibrary{positioning,arrows.meta}
\begin{tikzpicture}[
    stage/.style={draw=black!55, rounded corners=3pt, fill=black!3, align=center, minimum height=19mm,
                  text width=23mm, inner sep=3pt, font=\footnotesize},
    model/.style={stage, fill=orange!15, draw=orange!75!black},
    file/.style={font=\scriptsize\ttfamily, text=black!60, align=center, text width=24mm, anchor=south},
    >={Stealth[length=2mm]}
]
    \node[stage] (s1) {\textbf{text8}\\first $N$ tokens};
    \node[stage, right=5mm of s1] (s2) {\textbf{Vocabulary}\\count $\geq 5$,\\subsampling};
    \node[stage, right=5mm of s2] (s3) {\textbf{Training pairs}\\dynamic window\\$W \sim U\{1..5\}$};
    \node[model, right=5mm of s3] (s4) {\textbf{Negative}\\\textbf{sampling}\\$K = 5$, $f(w)^{3/4}$};
    \node[model, right=5mm of s4] (s5) {\textbf{SGNS step}\\forward, gradients,\\SGD update};
    \node[stage, right=5mm of s5] (s6) {\textbf{Embeddings}\\$W_{\text{in}} \in \mathbb{R}^{V \times d}$};
    \node[stage, right=5mm of s6] (s7) {\textbf{Evaluation}\\neighbours, cosine,\\analogy, PCA};

    \foreach \a/\b in {s1/s2, s2/s3, s3/s4, s4/s5, s5/s6, s6/s7}
        \draw[->, black!70] (\a) -- (\b);
    \draw[->, orange!75!black] (s5.south) -- ++(0,-3mm) -| node[pos=0.25, below, font=\scriptsize, text=orange!60!black] {repeat for every pair, epochs} (s3.south);

    \node[file, above=1mm of s1] {dataset/\\text8.py};
    \node[file, above=1mm of s2] {dataset/\\vocabulary.py};
    \node[file, above=1mm of s3] {model/\\word2vec.py};
    \node[file, above=1mm of s4] {\\utils.py};
    \node[file, above=1mm of s5] {model/\\word2vec.py};
    \node[file, above=1mm of s6] {training/\\train.py};
    \node[file, above=1mm of s7] {evaluation/\\evaluate.py};
\end{tikzpicture}
}
    \caption{Training and evaluation pipeline. File names are relative to \texttt{src/}. The orange stages run once per training pair.}
    \label{fig:pipeline}
\end{figure}

\subsection{Dataset}

We use text8 \citep{mahoney_text8}, a cleaned extract of English Wikipedia distributed for compression benchmarks: lower-case letters \texttt{a}--\texttt{z} separated by single spaces, about 17 million tokens in one line, with punctuation removed and digits spelled out as words. Tokenisation is therefore a whitespace split. To keep training short, a run uses only the first $N$ tokens of the file (\DemoTokens{} or \FinalTokens{} in our experiments), so it sees a contiguous prefix of the corpus rather than a random sample of articles. The corpus is downloaded automatically on first use.

\subsection{Vocabulary and subsampling}

The vocabulary contains every word that occurs at least 5 times in the $N$ tokens, indexed by decreasing frequency; other words are dropped from the token sequence. Frequent words are then subsampled as in Section~\ref{sec:background-ns}: each occurrence of $w$ is kept with probability $\min\!\left(1, \sqrt{t / f(w)}\right)$, $t = 10^{-3}$, where $f(w)$ is the relative frequency of $w$ in the vocabulary. Dropping rare words and subsampling together leave about two thirds of the tokens (Table~\ref{tab:runs}). It is applied once before training, so all epochs see the same subsampled sequence rather than a new draw per pass.

\subsection{Training pairs}

Pairs are generated on the fly. For each position $i$ of the subsampled sequence, a window size $W$ is drawn uniformly from $\{1, \dots, 5\}$ and the centre word is paired with the $W$ words on each side. Drawing $W$ at random weights nearby words more, since the immediate neighbours are inside the window for every $W$ while words five positions away are included only when $W = 5$. A centre word yields $W + 1$ context words on average, i.e.\ 6 with a maximum window of 5, so one epoch has about $6N'$ training pairs for $N'$ tokens after subsampling.

\subsection{Model, loss and gradients}

$W_{\text{in}}$ is initialised uniformly in $[-0.5/d, 0.5/d]$ and $W_{\text{out}}$ with zeros. For a centre word $c$, a positive context word $o$ and negative words $n_1, \dots, n_K$, let $s_o = \mathbf{v}'_o \cdot \mathbf{v}_c$ and $s_k = \mathbf{v}'_{n_k} \cdot \mathbf{v}_c$. The code minimises the negative of objective~\eqref{eq:sgns},
\begin{equation}
    L = -\log \sigma(s_o) - \sum_{k=1}^{K} \log\bigl(1 - \sigma(s_k)\bigr),
\end{equation}
which is the same quantity because $1 - \sigma(x) = \sigma(-x)$. Using $\sigma'(x) = \sigma(x)(1 - \sigma(x))$, the gradients are
\begin{align}
    \frac{\partial L}{\partial \mathbf{v}_c} &= \bigl(\sigma(s_o) - 1\bigr)\, \mathbf{v}'_o + \sum_{k=1}^{K} \sigma(s_k)\, \mathbf{v}'_{n_k}, \\
    \frac{\partial L}{\partial \mathbf{v}'_o} &= \bigl(\sigma(s_o) - 1\bigr)\, \mathbf{v}_c,
    \qquad
    \frac{\partial L}{\partial \mathbf{v}'_{n_k}} = \sigma(s_k)\, \mathbf{v}_c .
\end{align}
Only $K + 2$ rows of the two matrices change per step. Listing~\ref{lst:step} shows the corresponding training step from the original repository, without its comments. \texttt{np.add.at} accumulates the updates correctly when the same negative word is drawn twice.

\begin{lstlisting}[caption={One SGNS training step (\texttt{MyWord2Vec.step} in \texttt{src/model/word2vec.py}).}, label={lst:step}, float=t]
v_c = self.W_in[centre_idx]                  # (D)
v_o = self.W_out[pos_idx]                    # (D)
V_n = self.W_out[neg_idxs]                   # (K, D)

sig_o = sigmoid(v_c @ v_o)                   # P(positive | centre)
sig_k = sigmoid(V_n @ v_c)                   # (K)
loss = -np.log(sig_o + eps) - np.sum(np.log(1.0 - sig_k + eps))

grad_s_o = sig_o - 1.0                       # dL/ds_o
grad_s_k = sig_k                             # dL/ds_k
grad_vc = grad_s_o * v_o + grad_s_k @ V_n
grad_vo = grad_s_o * v_c
grad_Vn = np.outer(grad_s_k, v_c)            # (K, D)

self.W_in[centre_idx] -= lr * grad_vc        # SGD updates
self.W_out[pos_idx]   -= lr * grad_vo
np.add.at(self.W_out, neg_idxs, -lr * grad_Vn)
\end{lstlisting}

\subsection{Negative sampling in practice}
\label{sec:method-negatives}

The noise distribution is the unigram distribution of the vocabulary raised to the power $3/4$ and renormalised. Each positive pair receives $K = 5$ negatives; a negative that coincides with the centre or the context word is re-drawn.

The original repository drew each set of negatives with \texttt{np.random.choice(V, size=K, p=noise\_dist)}. This call rebuilds the cumulative distribution over the whole vocabulary every time, so its cost grows with $V$ and is paid once per training pair. Profiling showed that it took most of the training time: throughput fell from about 10,900 steps per second with \DemoVocab{} words to about 5,200 with \FinalVocab{} words, while the gradient step itself took only a small share. We replaced it with a sampler that draws one million i.i.d.\ negatives at a time from the same distribution and hands them out in order. The samples follow the same distribution: over two million draws, their total variation distance to $P_n$ was 0.041, identical to that of direct sampling, which is pure sampling noise. Throughput became 41,000--46,000 steps per second in our runs, almost independent of vocabulary size (Table~\ref{tab:runs}), and the \DemoTokensShort-token, three-epoch configuration took 251 seconds instead of 1,023 seconds with the original code. Both versions reached the same per-epoch losses (Section~\ref{sec:results-loss}), confirming that training behaves the same.

\subsection{Optimisation}

Training uses plain stochastic gradient descent on one pair at a time. The learning rate decays linearly from $\eta_0 = 0.025$ towards $10^{-4}$ over the estimated number of steps $\text{epochs} \times N' \times 6$:
\begin{equation}
    \eta_t = \max\!\left(10^{-4}, \; \eta_0 \left(1 - t / T\right)\right).
\end{equation}
After training, the rows of $W_{\text{in}}$ serve as the word embeddings; $W_{\text{out}}$ is used only during training.
